\subsection{Experimental setup}\label{sec:setup}

This subsection outlines the experimental framework designed to evaluate the classification performance, dimensionality reduction capability, computational efficiency, and hyperparameter robustness of KDS-EL.

\subsubsection{High-dimensional benchmark datasets}

To validate the performance and generalization ability of KDS-EL across diverse problem domains, a total of 20 benchmark datasets widely adopted in the evolutionary feature selection literature~\cite{mel2024,wangjia2022,wei2025} are selected. These datasets encompass diverse categories and sample sizes across biomedical gene-expression microarrays, text categorization corpora, image recognition, speech processing, and synthetic benchmarks. Such diverse data distributions and dimensional complexities effectively evaluate the generalization capability and feature reduction efficacy of feature selection algorithms under challenging conditions. The detailed characteristics of the 20 benchmark datasets are summarized in Table~\ref{T2}.

\subsubsection{Comparison algorithms}

We select two primary categories of comparison algorithms for analysis with the proposed KDS-EL framework, as summarized in Table~\ref{tab:comparison_literature}. The first category comprises specialized single-objective evolutionary algorithms in the field of high-dimensional feature selection, which includes: FWPSO~\cite{wangjia2022}, PSO-CSM~\cite{wei2025}, SaWDE~\cite{wang2022}, IGPSO~\cite{xuezhang2024}, and BPSO~\cite{kennedy1997}. The relevant parameter settings are shown in Table~\ref{T3}.

Another category of representative algorithms encompasses multi-objective and multi-task feature selection approaches, including NSGA-II-FS~\cite{deb2002}, MOEA/D-FS~\cite{zhang2007}, MOPSO-FS~\cite{coello2004}, fMOPSO-FS~\cite{han2024}, W-MOPSO-FS~\cite{xue2012mopso}, and MEL~\cite{mel2024}. In addition, two constituent variants of KDS-EL are investigated in Section~\ref{sec:improvement} to evaluate individual component contributions. Together with established single-objective and multi-objective competitors, this 11-algorithm comparison suite systematically spans classical search paradigms and state-of-the-art developments published between 2022 and 2025, ensuring a balanced, rigorous, and reproducible empirical evaluation.

\subsubsection{Evaluation metrics and statistical test protocols}

Algorithm performance is evaluated across three primary quantitative criteria: mean classification accuracy alongside standard deviation, the number of selected features alongside the feature selection ratio ($\mathrm{FSR} = \frac{\bar{d}}{D} \times 100\%$), and computational runtime in seconds. The $k$-nearest neighbor classifier with $k=3$ serves as the base evaluation model due to its non-parametric nature, parameter-free reproducibility, and high sensitivity to candidate feature subset discriminability. Following the standard benchmarking protocol established in high-dimensional evolutionary feature selection literature~\cite{mel2024,wangjia2022,wei2025}, candidate feature subsets are evaluated using 5-fold stratified cross-validation as the internal wrapper fitness surrogate across all problem instances. Under this standardized formulation, each benchmark dataset $\mathcal{D}$ is treated as a self-contained combinatorial optimization instance; the 5-fold cross-validation classification accuracy guides the evolutionary search trajectory and evaluates the discriminative quality of the optimal feature subsets under identical computational budgets across all competing methods. To guarantee a rigorous and methodologically equitable comparison between single-objective wrappers and multi-objective evolutionary algorithms that yield an entire Pareto non-dominated front $\mathcal{PF}$, a standardized post-Pareto decision-making protocol is enforced. In accordance with established multi-objective feature selection benchmarks~\cite{xue2012mopso,han2024}, the final operational solution $\mathbf{x}^{\star} \in \mathcal{PF}$ is uniquely extracted as the non-dominated candidate that achieves the minimal classification error on the cross-validation folds, with ties broken by smaller subset cardinality:
\begin{equation}
\mathbf{x}^{\star} = \arg\min_{\mathbf{x} \in \mathcal{PF}} \Bigl(\mathrm{ErrorRate}(\mathbf{x}), \; \|\mathbf{x}\|_0\Bigr).
\label{eq:pareto_selection}
\end{equation}
This operational protocol ensures that all competing paradigms are systematically evaluated at their maximum predictive accuracy.

To guarantee experimental fairness and statistical reliability, every algorithm is executed with a population size of $N=20$ candidate solutions over $T_{\max}=100$ generations across 10 independent runs with distinct random seeds. In ultra-high-dimensional wrapper feature selection, each fitness evaluation requires training and validating a classifier across thousands or tens of thousands of candidate dimensions (e.g., up to 54,675 dimensions in \textit{T1D}), imposing an exceptionally intensive computational burden of $N \times T_{\max} = 2{,}000$ cross-validation evaluations per single run. Executing 10 independent runs per dataset corresponds to 50 distinct stratified cross-validation folds, which provides an adequate sample size to capture classifier generalization variance and algorithmic stability while remaining computationally tractable across an extensive 14-algorithm comparison suite. This 10-run protocol strictly aligns with the benchmark evaluation standard established in recent top-tier literature on ultra-high-dimensional feature selection, most notably the multi-task evolutionary learning framework MEL published in \textit{IEEE Transactions on Knowledge and Data Engineering}~\cite{mel2024}, thereby guaranteeing an equitable and reproducible empirical comparison.

Furthermore, to determine whether the empirical advantages of KDS-EL over competing paradigms are statistically significant rather than stochastic artifacts, we strictly follow the multi-problem non-parametric statistical testing framework recommended by Dem{\v{s}}ar~\cite{demsar2006}. In this protocol, statistical power is governed by the number of independent benchmark problems ($N_{\mathrm{datasets}} = 20$), with each problem represented by the paired summary performance across the independent runs. Specifically, the non-parametric Friedman test~\cite{friedman1940} is utilized to establish global omnibus algorithm rankings across all 20 benchmark datasets. Pairwise Wilcoxon signed-rank tests at the $\alpha=0.05$ significance level assess directional superiority between KDS-EL and each competitor, with the Holm--Bonferroni step-down procedure~\cite{holm1979} applied to rigorously control the family-wise error rate across multiple hypothesis comparisons.
